\section{Translation and Evaluation Pipeline}

The translation and evaluation process follows a consistent pipeline for both
the base and fine-tuned models. Each model receives Java input code from the
translation subset of CodeEditorBench and generates the corresponding C++
translation. The input examples are formatted using the model-specific prompt
format described in the experimental setup.

After preprocessing, the models generate C++ translations from the Java input
code samples contained in the CodeEditorBench dataset. The generated outputs
are then post-processed to remove markdown code fences, explanations, and any
additional text outside the generated C++ code.

The generated translations are evaluated using both similarity-based and
execution-based approaches. First, BLEU score is computed by comparing the
generated C++ output against the reference \texttt{target\_code} solution from
the evaluation dataset.

Execution-based evaluation is then performed using the CodeEditorBench
framework.

For each model, both BLEU scores and execution-based metrics are collected. The
results are compared across the base and fine-tuned variants of StarCoder2 and
Qwen in order to analyze the effect of fine-tuning.